% =============================================================================
%  Notas del Profesor — Sesión 19: Proyectos de múltiples archivos
%  Programación Avanzada — Universidad Panamericana
%  Compilar: pdflatex sesion19_proyectos_multiarchivo_profesor.tex
% =============================================================================
\documentclass[12pt, oneside]{article}
\usepackage[profesor]{progavanzada}
\usepackage{array}

\begin{document}

\portada{Sesión 19}{Proyectos de múltiples archivos}{2026}

\tableofcontents
\newpage

% =============================================================================
\section{Objetivos de la sesión}
% =============================================================================

Al finalizar la sesión el alumno será capaz de:

\begin{enumerate}
  \item Explicar la diferencia entre declaración y definición en C++ y por
        qué los headers contienen solo declaraciones.
  \item Escribir un archivo \texttt{.h} correcto con guard de inclusión.
  \item Describir los dos pasos de compilación (\texttt{.cpp} → \texttt{.o}
        → ejecutable) y el rol del enlazador.
  \item Escribir un \texttt{Makefile} básico partiendo de la línea de
        compilación directa y evolucionarlo con variables.
  \item Crear un \texttt{CMakeLists.txt} mínimo y seguir el flujo
        \textit{out-of-source build}.
  \item Leer un proyecto de varios módulos e identificar el grafo de
        dependencias entre archivos.
\end{enumerate}

\begin{notaprofesor}[Duración estimada]
  90 minutos.  Distribución sugerida:\\[2pt]
  \begin{tabular}{rl}
    10 min & gancho histórico + preguntas de arranque \\
    15 min & headers: declaración vs.\ definición, guards \\
    20 min & demo en vivo: compilar por etapas con \texttt{g++} \\
           & + la línea que no para de crecer \\
    25 min & Make: del Makefile mínimo a las variables \\
    15 min & CMake: \texttt{CMakeLists.txt} y flujo de trabajo \\
     5 min & proyecto de personas + asignación de ejercicios \\
  \end{tabular}\\[4pt]
  Los alumnos siguen la demo en el cuaderno Jupyter
  (\texttt{sesiones/sesion19\_proyectos\_multiarchivo.ipynb}).
  Los ejercicios se terminan en casa.
\end{notaprofesor}

% =============================================================================
\section{Prerrequisitos}
% =============================================================================

\begin{itemize}
  \item Compilar y ejecutar un programa de un solo archivo con \texttt{g++}
        (sesiones 1--2).
  \item Funciones: declaración y llamada (sesión 4).
  \item Apuntadores y arrays como parámetros (sesión 14).
  \item Familiaridad básica con terminal: \texttt{ls}, \texttt{cd},
        \texttt{cat}, redirección.
\end{itemize}

\begin{notaprofesor}
  Esta sesión no requiere recursión ni conocimiento de STL.  Es intencionalmente
  accesible: el arco conceptual va de ``un archivo'' a ``varios archivos'', y todo
  el herramental (headers, \texttt{.o}, Make, CMake) es motivado por ese problema.
  \par\smallskip
  Señal de alerta en la demo: si un alumno no recuerda cómo compilar con
  \texttt{g++}, detenerse brevemente en los flags más comunes (\texttt{-o},
  \texttt{-Wall}, \texttt{-std=c++17}) antes de agregar \texttt{-c}.
\end{notaprofesor}

% =============================================================================
\section{Arranque y preguntas de sondeo}
% =============================================================================

\begin{notaprofesor}[Preguntas de arranque (10 min)]
  \textbf{1.\ ``¿Cuántos archivos tiene el proyecto más grande que han escrito?''}\\
  Respuesta esperada: uno.  Si alguien dice más, pedirle que explique cómo
  lo organizó.  Seguir con: ``¿cómo saben en cuál archivo está la función
  que necesitan?''

  \medskip
  \textbf{2.\ ``Si tienen un archivo de 2\,000 líneas y cambian una función
  de 3 líneas, ¿qué tiene que hacer la computadora?''}\\
  Respuesta esperada: recompilar todo.  Seguir con: ``¿y si pudiéramos
  recompilar solo esa función?'' — eso es lo que aprenderemos hoy.

  \medskip
  \textbf{3.\ ``¿Cómo creen que funciona el kernel de Linux, que tiene más
  de 30 millones de líneas?''}\\
  Respuesta esperada: muchos archivos.  Dato de gancho: el kernel tiene
  $\approx$70\,000 archivos \texttt{.c} y su \texttt{Makefile} tiene miles
  de líneas.  Sin herramientas de gestión de compilación sería imposible
  de mantener.
\end{notaprofesor}

% =============================================================================
\section{Declaración versus definición}
% =============================================================================

\subsection{Explicación central}

La distinción es sencilla pero se confunde con frecuencia:

\begin{itemize}
  \item \textbf{Declaración:} le dice al compilador que algo \textit{existe}
        y cuál es su forma (tipo de retorno, parámetros).  No genera código
        ejecutable.  Se puede repetir cuantas veces sea necesario.
  \item \textbf{Definición:} proporciona el cuerpo.  Solo puede aparecer
        \textbf{una vez} en todo el programa (\textbf{ODR}: \textit{One Definition Rule}).
\end{itemize}

\begin{notaprofesor}[Analogía recomendada: el índice del libro]
  ``Una \textit{declaración} es como el índice de un libro: te dice que
  existe el capítulo 5 y en qué página está.  La \textit{definición} es
  el capítulo completo.  Puedes imprimir el índice mil veces (en cada
  archivo que lo necesite), pero el capítulo completo solo existe una vez.''
  \par\smallskip
  Esta analogía funciona bien para motivar el patrón
  \texttt{.h} = declaraciones, \texttt{.cpp} = definiciones.
\end{notaprofesor}

\subsection{Guards de inclusión}

\begin{notaprofesor}[Demo recomendada: error sin guard (5 min)]
  Crear un \texttt{funciones.h} sin guard que declare \texttt{void saludar(char*);}.
  En \texttt{main.cpp}, incluirlo dos veces:
\begin{verbatim}
#include "funciones.h"
#include "funciones.h"   // segunda inclusion deliberada
\end{verbatim}
  Con solo declaraciones no ocurre ningún error (las declaraciones se pueden
  repetir).  Ahora poner una definición en el header:
\begin{verbatim}
void saludar(char* n) { cout << "Hola " << n; }
\end{verbatim}
  Compilar y mostrar el error:
\begin{verbatim}
error: redefinition of 'void saludar(char*)'
\end{verbatim}
  Agregar el guard y demostrar que desaparece.  Luego quitar la definición del
  header (regla: solo declaraciones) y verificar que tampoco hay error sin guard.
  \par\smallskip
  \textbf{Por qué vale la pena hacer esto:} los alumnos recuerdan el guard
  porque lo ven solucionar un error concreto, no porque alguien les dijo que
  lo pusieran.
\end{notaprofesor}

\begin{notaprofesor}[FAQ: \texttt{\#pragma once} vs.\ guard manual]
  \textbf{¿Cuál es mejor?}\\
  \texttt{\#pragma once} es más simple y elimina el riesgo de escribir mal
  el nombre de la macro (error clásico: \texttt{FUNCIONES\_H} en el
  \texttt{\#ifndef} y \texttt{FUNCIONES\_HH} en el \texttt{\#define}).
  El guard manual es el estándar ISO histórico (aparece en todos los
  tutoriales y código existente).  Para este curso, ambas formas son
  aceptables.

  \medskip
  \textbf{¿Puede un \texttt{.h} contener definiciones?}\\
  Sí, en dos casos: funciones \texttt{inline} y plantillas (\textit{templates}).
  Ambos tienen reglas especiales en la ODR que permiten múltiples
  definiciones idénticas.  Para este curso, la regla simple es:
  solo declaraciones en el \texttt{.h}.
\end{notaprofesor}

% =============================================================================
\section{Compilación por etapas}
% =============================================================================

\subsection{Demo en vivo: los dos pasos}

\begin{notaprofesor}[Guión de la demo (15 min)]
  Trabaja con los archivos de \texttt{materialInicial/sesion19/}
  (\texttt{main.cpp}, \texttt{funciones.h}, \texttt{funciones.cpp}).

  \begin{enumerate}
    \item Mostrar los tres archivos con \texttt{cat} o en el editor.
    \item Compilar solo \texttt{funciones.cpp}:
\begin{verbatim}
g++ -c funciones.cpp -o funciones.o
ls -lh funciones.o
\end{verbatim}
    \item Intentar ejecutar el \texttt{.o} directamente:
\begin{verbatim}
./funciones.o
\end{verbatim}
          Da error — no es un ejecutable.  Preguntar: ``¿qué falta?''
          El punto de entrada (\texttt{main}).  Eso es exactamente lo
          que aporta el enlazador.
    \item Opcional: \texttt{nm --demangle funciones.o} para mostrar los
          símbolos definidos (\texttt{T}) en el objeto.
    \item Compilar \texttt{main.cpp}:
\begin{verbatim}
g++ -c main.cpp -o main.o
nm --demangle main.o | grep saludar
\end{verbatim}
          Señalar que \texttt{saludar} aparece como \texttt{U}
          (\textit{undefined}) en \texttt{main.o}: está declarada pero
          no definida ahí.
    \item Enlazar:
\begin{verbatim}
g++ main.o funciones.o -o programa
./programa
\end{verbatim}
    \item Modificar \texttt{funciones.cpp} (cambiar el mensaje).
          Mostrar que solo hay que recompilar ese archivo:
\begin{verbatim}
g++ -c funciones.cpp -o funciones.o
g++ main.o funciones.o -o programa
./programa
\end{verbatim}
          \texttt{main.o} no se volvió a compilar.  Preguntar:
          ``¿por qué no necesitamos recompilar \texttt{main.o}?''
  \end{enumerate}
\end{notaprofesor}

\subsection{La línea que no para de crecer}

\begin{notaprofesor}[Transición clave: motivar Make (5 min)]
  Después de la demo, mostrar cómo se ve la línea de compilación para el
  proyecto de personas (6 módulos):
\begin{verbatim}
g++ -std=c++17 -Wall -Wextra -o programa \
    main.cpp auxiliar.cpp obtener_datos.cpp \
    mate.cpp buscar.cpp
\end{verbatim}
  Preguntar: ``¿Qué pasa cuando el proyecto tenga 20 archivos?  ¿O cuando
  necesitemos agregar \texttt{-O2} para optimización?  ¿O linkear una
  librería externa?''  Mostrar la versión inflada:
\begin{verbatim}
g++ -std=c++17 -Wall -Wextra -O2 -I./include \
    -I./lib/json/include \
    main.cpp auxiliar.cpp obtener_datos.cpp    \
    mate.cpp buscar.cpp graficas.cpp reportes.cpp \
    -L./lib -ljson -lm -o programa
\end{verbatim}
  Pausa.  ``Hay una solución mejor.  Se llama \texttt{Makefile}.''
  \par\smallskip
  \textbf{Por qué funciona esta transición:} el alumno ya vivió la incomodidad
  de la línea larga durante la demo.  La motivación es visceral, no abstracta.
\end{notaprofesor}

\begin{notaprofesor}[FAQ: compilación]
  \textbf{¿Qué es un ``error de enlace'' (\textit{linker error})?}\\
  Ocurre cuando el linker no encuentra la definición de algo declarado.
  El mensaje típico es:
\begin{verbatim}
undefined reference to `saludar(char*)'
\end{verbatim}
  Causa más común: olvidar incluir un \texttt{.cpp} en el comando de enlace.

  \medskip
  \textbf{¿Por qué el \texttt{.h} no se compila directamente?}\\
  El \texttt{.h} no contiene definiciones; el preprocesador lo copia
  textualmente en cada \texttt{.cpp} que lo incluye.  Técnicamente puedes
  hacer \texttt{g++ -c funciones.h}, pero genera un objeto vacío o inútil.

  \medskip
  \textbf{¿Qué significa el flag \texttt{-I./include}?}\\
  Le dice al compilador dónde buscar headers adicionales.  Sin ese flag,
  \texttt{\#include "json.hpp"} fallaría si el archivo no está en el
  directorio actual.
\end{notaprofesor}

% =============================================================================
\section{Make: del Makefile mínimo a las variables}
% =============================================================================

\subsection{Secuencia de escritura en vivo}

\begin{notaprofesor}[Guión de la demo de Make (25 min)]
  Escribe el Makefile en vivo, paso a paso.  No muestres el resultado final
  de entrada — la evolución gradual es el punto pedagógico central de esta sesión.

  \medskip
  \textbf{Paso 1 — Solo la línea de compilación (5 min)}

  Crear \texttt{Makefile} con el editor y escribir:
\begin{verbatim}
programa:
	g++ -std=c++17 -Wall -Wextra -o programa \
	    main.cpp auxiliar.cpp obtener_datos.cpp \
	    mate.cpp buscar.cpp
\end{verbatim}
  Ejecutar \texttt{make}.  Preguntar: ``¿Qué ganamos respecto al comando
  directo?''  Respuesta: no tenemos que \textit{recordar} ni \textit{teclear}
  la línea — solo \texttt{make}.
  \par\smallskip
  Señalar el Tab antes del comando.  Borrar el Tab, reemplazar con espacios
  y ejecutar de nuevo para provocar el error:
\begin{verbatim}
Makefile:2: *** missing separator.  Stop.
\end{verbatim}
  Restaurar el Tab.

  \medskip
  \textbf{Paso 2 — Variable para el compilador (3 min)}

  Agregar \texttt{CXX = g++} arriba y sustituir en el comando.  Preguntar:
  ``¿Por qué es útil poner el compilador en una variable?''  Si alguien
  trabaja con \texttt{clang++} o necesita cross-compilar, solo cambia una
  línea.

  \medskip
  \textbf{Paso 3 — Variable para los flags (3 min)}

  Agregar \texttt{CXXFLAGS = -std=c++17 -Wall -Wextra}.  Enfatizar: cuando
  el proyecto crezca y se necesite agregar \texttt{-O2} o \texttt{-g},
  se edita solo \texttt{CXXFLAGS}, no cada regla.

  \medskip
  \textbf{Paso 4 — Variables para archivos y target (4 min)}

  Agregar \texttt{TARGET} y \texttt{SRCS}:
\begin{verbatim}
CXX      = g++
CXXFLAGS = -std=c++17 -Wall -Wextra
TARGET   = programa
SRCS     = main.cpp auxiliar.cpp obtener_datos.cpp \
           mate.cpp buscar.cpp

$(TARGET): $(SRCS)
	$(CXX) $(CXXFLAGS) $(SRCS) -o $(TARGET)
\end{verbatim}
  Preguntar: ``Ahora quiero agregar \texttt{graficas.cpp}.  ¿Qué línea
  toco?''  Solo \texttt{SRCS}.

  \medskip
  \textbf{Paso 5 — Introducir OBJS y recompilación selectiva (10 min)}

  Ejecutar \texttt{make} dos veces seguidas para mostrar que la segunda
  vez recompila todo aunque nada cambió.  Preguntar: ``¿Cómo podría Make
  saber que ya está compilado?''  Llevar hacia la idea de fechas y archivos
  intermedios.

  Evolucionar a la versión con \texttt{OBJS}:
\begin{verbatim}
OBJS = main.o auxiliar.o obtener_datos.o mate.o buscar.o

$(TARGET): $(OBJS)
	$(CXX) $(OBJS) -o $(TARGET)

%.o: %.cpp
	$(CXX) $(CXXFLAGS) -c $< -o $@

clean:
	rm -f $(OBJS) $(TARGET)
\end{verbatim}
  Ejecutar \texttt{make}.  Luego:
\begin{verbatim}
touch auxiliar.cpp    # simula una edición en auxiliar.cpp
make
\end{verbatim}
  Mostrar que solo \texttt{auxiliar.o} se recompila.  Luego:
\begin{verbatim}
touch persona.h       # simula un cambio en el header compartido
make
\end{verbatim}
  Mostrar que Make \textit{no} detecta esto sin dependencias explícitas
  de headers — solo reenlaza el ejecutable.  Decir: ``¿Es correcto?  No.
  Pero Fix completo con dependencias de headers es para el ejercicio.''
  (El alumno lo resuelve en el ejercicio 2.)
  \par\smallskip
  Terminar con \texttt{make clean} y verificar que desaparecen los artefactos.
\end{notaprofesor}

\subsection{Makefile completo de referencia}

\begin{verbatim}
CXX      = g++
CXXFLAGS = -std=c++17 -Wall -Wextra
TARGET   = programa
SRCS     = main.cpp auxiliar.cpp obtener_datos.cpp mate.cpp buscar.cpp
OBJS     = $(SRCS:.cpp=.o)

$(TARGET): $(OBJS)
	$(CXX) $(OBJS) -o $(TARGET)

%.o: %.cpp
	$(CXX) $(CXXFLAGS) -c $< -o $@

clean:
	rm -f $(OBJS) $(TARGET)
\end{verbatim}

\begin{notaprofesor}[FAQ: Make]
  \textbf{Error ``missing separator'':}\\
  El comando empieza con espacios, no con Tab.  En VS Code se ve igual
  visualmente; decir a los alumnos que abran el archivo en \texttt{cat -A}
  para ver los caracteres ocultos (\texttt{\^{}I} = Tab, espacios = espacios).

  \medskip
  \textbf{¿Por qué no usar simplemente \texttt{g++ *.cpp} siempre?}\\
  Funciona para compilar, pero recompila \textit{todo} en cada invocación.
  En proyectos de 100+ archivos con 5+ minutos de compilación, es inaceptable.

  \medskip
  \textbf{¿Qué significa \texttt{\$(SRCS:.cpp=.o)}?}\\
  Es una sustitución de texto en Make: reemplaza la extensión \texttt{.cpp}
  por \texttt{.o} en cada elemento de \texttt{SRCS}.  Equivale a escribir
  \texttt{OBJS = main.o auxiliar.o ...} a mano.

  \medskip
  \textbf{¿Cómo agrego dependencias de headers?}\\
  Declarando el \texttt{.o} como objetivo con sus \texttt{.h} como
  dependencias:
\begin{verbatim}
auxiliar.o: auxiliar.cpp auxiliar.h persona.h
\end{verbatim}
  Con esto, si \texttt{persona.h} cambia, Make recompila \texttt{auxiliar.o}.
  Sin esta línea, Make no lo sabe.  (Esto es el ejercicio 2.)

  \medskip
  \textbf{¿Hace Make algo con \texttt{clean}?}\\
  \texttt{clean} no es una regla especial de Make — es solo un objetivo
  sin archivo correspondiente.  Make lo ejecuta cuando escribes
  \texttt{make clean}.  La convención del nombre es universal.
\end{notaprofesor}

% =============================================================================
\section{CMake}
% =============================================================================

\begin{notaprofesor}[Demo de CMake (15 min)]
  \begin{enumerate}
    \item Crear \texttt{CMakeLists.txt} en el directorio del proyecto:
\begin{verbatim}
cmake_minimum_required(VERSION 3.10)
project(MiProyecto VERSION 1.0)
set(CMAKE_CXX_STANDARD 17)
set(CMAKE_CXX_STANDARD_REQUIRED True)
add_executable(programa
    main.cpp
    auxiliar.cpp
    obtener_datos.cpp
    mate.cpp
    buscar.cpp
)
\end{verbatim}
    \item Crear el directorio de build y configurar:
\begin{verbatim}
mkdir build && cd build
cmake ..
\end{verbatim}
          Señalar los archivos generados: \texttt{Makefile},
          \texttt{CMakeCache.txt}, \texttt{CMakeFiles/}.
    \item Compilar y ejecutar:
\begin{verbatim}
make
./programa
\end{verbatim}
    \item Regresar al directorio fuente (\texttt{cd ..}) y mostrar que
          está limpio: solo \texttt{CMakeLists.txt}, los \texttt{.cpp}
          y los \texttt{.h}.  Ningún \texttt{.o} ni \texttt{Makefile}
          generado ensucia el directorio de código.
    \item Preguntar: ``¿qué poner en \texttt{.gitignore}?''
          Respuesta: \texttt{build/}.
  \end{enumerate}
\end{notaprofesor}

\begin{notaprofesor}[FAQ: CMake]
  \textbf{¿Para qué sirve \texttt{cmake\_minimum\_required}?}\\
  Protege contra versiones antiguas de CMake que no soporten las funciones
  usadas en el archivo.  Buena práctica: poner la versión mínima real que
  necesita el proyecto.

  \medskip
  \textbf{¿Ninja vs.\ Make en CMake?}\\
  \texttt{cmake -GNinja ..} genera archivos Ninja en lugar de Makefile.
  Ninja es más rápido en proyectos grandes (mejor paralelismo).
  Para este curso, Make es suficiente.

  \medskip
  \textbf{¿Cómo agrego flags de compilación en CMake?}\\
\begin{verbatim}
target_compile_options(programa PRIVATE -Wall -Wextra)
\end{verbatim}
  O globalmente con \texttt{add\_compile\_options(-Wall)}.

  \medskip
  \textbf{``Prefiero no usar \texttt{mkdir build / cd build}, ¿hay otra forma?''}\\
  Sí: \texttt{cmake -B build} desde el directorio raíz genera el build
  en \texttt{build/} sin necesidad de \texttt{cd}.  Luego
  \texttt{cmake --build build} compila sin entrar al directorio.
  Más cómodo en proyectos donde se abre el editor en la raíz.

  \medskip
  \textbf{¿Cuándo usar Make y cuándo CMake?}\\
  Make: proyecto propio que solo se usará en Linux/Mac, equipo pequeño,
  simplicidad máxima.  CMake: proyecto que debe funcionar en Windows,
  que se integrará con un IDE (CLion, VS, etc.), o que crecerá a más de
  un par de directorios.  La industria usa CMake casi universalmente en
  proyectos C++ modernos.
\end{notaprofesor}

% =============================================================================
\section{Proyecto ejemplo: sistema de personas}
% =============================================================================

\begin{notaprofesor}[Advertencia sobre el código de proyecto1]
  El código en \texttt{materialInicial/sesion19/proyecto1/} tiene inconsistencias
  que conviene conocer antes de la sesión.  El objetivo es usar el proyecto para
  \textit{leer} el grafo de dependencias, no para compilarlo.

  \begin{itemize}
    \item \textbf{Firmas inconsistentes en \texttt{buscar.h} vs.\ \texttt{buscar.cpp}:}
          el header declara tres parámetros en algunas funciones pero la
          implementación usa dos.  Código refactorizado a medias.
    \item \textbf{Funciones duplicadas:} \texttt{buscar.cpp} contiene
          implementaciones de \texttt{promedio\_edades} y
          \texttt{promedio\_salarios} que ya están en \texttt{mate.cpp}.
          Si se enlazan juntos, falla con error de ODR.
    \item \textbf{Include incorrecto en \texttt{main.cpp}:}
          incluye \texttt{"obtener.h"} pero el archivo se llama
          \texttt{obtener\_datos.h}.
  \end{itemize}

  \textbf{Sugerencia:} mostrar el grafo de inclusiones en el pizarrón
  y señalar por qué \texttt{persona.h} está en el centro.  Si se quiere
  compilar en clase, corregir los errores antes.
\end{notaprofesor}

% =============================================================================
\section{Errores frecuentes}
% =============================================================================

\begin{notaprofesor}[Errores más comunes de los alumnos]

  \textbf{1.\ Poner definiciones en el header.}\\
  Síntoma: el proyecto compila bien si solo \texttt{main.cpp} incluye el
  \texttt{.h}, pero falla al enlazar con un segundo \texttt{.cpp}.
  Error: \texttt{multiple definition of `f()'}.\\
  Remedio: mover la definición al \texttt{.cpp}; en el \texttt{.h}
  dejar solo la declaración.

  \vspace{0.5em}
  \textbf{2.\ Olvidar el Tab en el Makefile.}\\
  Error: \texttt{Makefile:N: *** missing separator.  Stop.}\\
  Remedio: el comando de la regla \textit{siempre} empieza con Tab.
  En VS Code, verificar que el archivo no esté en modo ``Indent with
  spaces''.

  \vspace{0.5em}
  \textbf{3.\ No incluir el \texttt{.h} propio en el \texttt{.cpp}.}\\
  Síntoma: el \texttt{.cpp} compila solo, pero al enlazar con \texttt{main.cpp}
  (que sí incluye el \texttt{.h}) el linker reporta que las firmas no coinciden.
  Remedio: el primer \texttt{\#include} en cualquier \texttt{.cpp} debe ser
  su propio \texttt{.h}.  Esto hace que el compilador verifique que la
  declaración y la definición son consistentes.

  \vspace{0.5em}
  \textbf{4.\ Usar \texttt{<>} para headers propios.}\\
  Síntoma: \texttt{fatal error: funciones.h: No such file or directory}.\\
  Remedio: headers del proyecto van con \texttt{""}, no con \texttt{<>}.

  \vspace{0.5em}
  \textbf{5.\ Olvidar agregar el \texttt{.cpp} al Makefile o CMakeLists.}\\
  Síntoma: el código compila pero hay errores de enlace para funciones del
  módulo nuevo.\\
  Remedio: todo \texttt{.cpp} que define funciones usadas por otros debe
  aparecer en la lista de fuentes del sistema de compilación.

\end{notaprofesor}

% =============================================================================
\section{Soluciones a los ejercicios}
% =============================================================================

\begin{notaprofesor}[Ejercicio 1 — Evolución del Makefile]
  El alumno debe mostrar cuatro versiones sucesivas del Makefile.

  \textbf{Versión 1 — Solo la línea de compilación:}
\begin{verbatim}
programa:
	g++ -std=c++17 -Wall -Wextra -o programa \
	    main.cpp funciones.cpp
\end{verbatim}

  \textbf{Versión 2 — Agrega \texttt{CXX}:}
\begin{verbatim}
CXX = g++

programa:
	$(CXX) -std=c++17 -Wall -Wextra -o programa \
	    main.cpp funciones.cpp
\end{verbatim}

  \textbf{Versión 3 — Agrega \texttt{CXXFLAGS} y \texttt{TARGET}:}
\begin{verbatim}
CXX      = g++
CXXFLAGS = -std=c++17 -Wall -Wextra
TARGET   = programa

$(TARGET):
	$(CXX) $(CXXFLAGS) -o $(TARGET) main.cpp funciones.cpp
\end{verbatim}

  \textbf{Versión 4 — Agrega \texttt{SRCS}:}
\begin{verbatim}
CXX      = g++
CXXFLAGS = -std=c++17 -Wall -Wextra
TARGET   = programa
SRCS     = main.cpp funciones.cpp

$(TARGET): $(SRCS)
	$(CXX) $(CXXFLAGS) $(SRCS) -o $(TARGET)
\end{verbatim}

  Criterio de evaluación: cada versión debe ser funcional (compilar el
  proyecto) y mostrar un avance claro respecto a la anterior.
\end{notaprofesor}

\begin{notaprofesor}[Ejercicio 2 — Observar recompilación]
  \textbf{Modificar \texttt{funciones.cpp}:}\\
  Make recompila \texttt{funciones.o} y reenlaza el ejecutable.
  \texttt{main.o} no se toca porque su fecha es más reciente que la de
  \texttt{main.cpp}.

  \medskip
  \textbf{Modificar \texttt{funciones.h}:}\\
  Con el Makefile de la sección 5 (que tiene la regla patrón pero
  \textit{no} declara dependencias de headers explícitas), Make \textbf{no}
  detecta el cambio en \texttt{funciones.h} porque ningún \texttt{.o}
  lo lista como dependencia.  Solo recompila si el \texttt{.o}
  correspondiente es más nuevo que el \texttt{.h}, lo cual no siempre ocurre.

  Para que funcione correctamente, hay que agregar:
\begin{verbatim}
main.o:      main.cpp funciones.h
funciones.o: funciones.cpp funciones.h
\end{verbatim}
  Ahora si \texttt{funciones.h} cambia, Make recompila \textbf{ambos} objetos.
\end{notaprofesor}

\begin{notaprofesor}[Ejercicio 3 — Doble definición]
  \textbf{Mismo archivo:}
\begin{verbatim}
void f() {}
void f() {}   // error: redefinition of 'void f()'
\end{verbatim}
  Error de \textit{compilación}.  El compilador lo detecta antes de enlazar.

  \medskip
  \textbf{Archivos distintos:} si \texttt{a.cpp} y \texttt{b.cpp} definen
  ambos \texttt{void f() \{\}}, cada archivo compila sin error,
  pero al enlazar:
\begin{verbatim}
error: multiple definition of `f()'
\end{verbatim}
  Error de \textit{enlace}.  El linker detecta la violación de la ODR.
  Este es el tipo de error que tendría el \texttt{proyecto1} al enlazar
  \texttt{mate.o} y \texttt{buscar.o} (duplican la definición de
  \texttt{promedio\_edades}).
\end{notaprofesor}

\begin{notaprofesor}[Ejercicio 4 — Módulo \texttt{estadisticas}]
  \textbf{estadisticas.h:}
\begin{verbatim}
#ifndef ESTADISTICAS_H
#define ESTADISTICAS_H
#include "persona.h"

double salario_maximo(Persona arr[], int n);

#endif
\end{verbatim}

  \textbf{estadisticas.cpp:}
\begin{verbatim}
#include "estadisticas.h"

double salario_maximo(Persona arr[], int n) {
    double max = arr[0].salario;
    for (int i = 1; i < n; i++)
        if (arr[i].salario > max) max = arr[i].salario;
    return max;
}
\end{verbatim}

  \textbf{Makefile — agregar a \texttt{SRCS}:}
\begin{verbatim}
SRCS = main.cpp auxiliar.cpp obtener_datos.cpp \
       mate.cpp buscar.cpp estadisticas.cpp
\end{verbatim}

  \textbf{CMakeLists.txt — agregar a \texttt{add\_executable}:}
\begin{verbatim}
add_executable(programa
    ...
    estadisticas.cpp
)
\end{verbatim}
\end{notaprofesor}

\begin{notaprofesor}[Ejercicio 5 — \texttt{make} vs.\ \texttt{make clean}]
  \textbf{\texttt{make}}: construye los objetivos cuyas dependencias
  son más recientes que ellos mismos, o los que no existen aún.

  \textbf{\texttt{make clean}}: ejecuta \texttt{rm -f \$(OBJS) \$(TARGET)}.
  Elimina todos los artefactos.  Útil cuando:
  \begin{itemize}
    \item Se cambia un flag de compilación (\texttt{CXXFLAGS}).  Make no
          detecta cambios en variables, solo en fechas de archivos.
    \item Se quiere verificar que el proyecto compila desde cero
          (``compilación limpia'').
    \item Los archivos \texttt{.o} son de una plataforma diferente
          (por ejemplo, se copió el proyecto de otro sistema).
  \end{itemize}
\end{notaprofesor}

% =============================================================================
\section{Conexiones con temas posteriores}
% =============================================================================

\begin{itemize}
  \item \textbf{Programación orientada a objetos} — Las clases en C++
    viven naturalmente en módulos: \texttt{Persona.h} / \texttt{Persona.cpp}.
    El patrón que ven hoy (declaración en \texttt{.h}, definición en
    \texttt{.cpp}) es exactamente el que usarán para cada clase del curso OOP.
  \item \textbf{Librerías estáticas y dinámicas} — Los archivos \texttt{.o}
    que se producen hoy son los bloques que se empaquetan en una librería
    (\texttt{.a} estática o \texttt{.so}/\texttt{.dll} dinámica).
  \item \textbf{Header-only libraries} — Muchas librerías modernas de C++
    (Eigen, nlohmann/json, stb) son solo un \texttt{.h}.  Explotan la
    excepción \texttt{inline}/templates de la ODR para no requerir un
    \texttt{.cpp}.  El alumno ya tiene los conceptos para entender por qué
    eso es posible.
  \item \textbf{Compilación incremental en CI/CD} — Los sistemas de
    integración continua (GitHub Actions, GitLab CI) usan Make o CMake
    para recompilar solo lo cambiado en cada PR.  El tiempo de pipeline
    se reduce por el mismo principio de fechas que vieron hoy.
\end{itemize}

% =============================================================================
\section{Rúbrica de ejercicios}
% =============================================================================

\begin{center}
\begin{tabular}{p{7cm}cc}
  \toprule
  \textbf{Criterio} & \textbf{Puntos} & \textbf{Evidencia} \\
  \midrule
  \multicolumn{3}{l}{\textit{Ejercicio 1 — Evolución del Makefile}} \\
  \quad Versión 1 funcional (solo línea)           & 5 & \texttt{make} produce ejecutable \\
  \quad Versiones 2--4 muestran evolución clara    & 5 & variables correctas, uso de \texttt{\$(VAR)} \\
  \midrule
  \multicolumn{3}{l}{\textit{Ejercicio 2 — Observar recompilación}} \\
  \quad Identifica correctamente qué recompila     & 5 & con y sin dependencia de header \\
  \quad Explica por qué \texttt{.h} no se detecta & 5 & dependencias explícitas son necesarias \\
  \midrule
  \multicolumn{3}{l}{\textit{Ejercicio 3 — Doble definición}} \\
  \quad Distingue error de compilación vs.\ enlace & 5 & mensaje de error correcto en cada caso \\
  \quad Explica con la ODR                         & 5 & referencia a la regla \\
  \midrule
  \multicolumn{3}{l}{\textit{Ejercicio 4 — Módulo estadisticas}} \\
  \quad Guard correcto en el header                & 4 & nombre consistente \\
  \quad \texttt{.cpp} incluye su propio \texttt{.h}& 3 & primera línea del \texttt{.cpp} \\
  \quad Agregado a Makefile y CMakeLists.txt       & 3 & en ambos lugares \\
  \midrule
  \multicolumn{3}{l}{\textit{Ejercicio 5 — make vs.\ make clean}} \\
  \quad Diferencia correcta                        & 5 & qué hace cada uno \\
  \quad Menciona el caso de cambio de flag         & 5 & Make no detecta cambios en variables \\
  \midrule
  \textbf{Total}                                   & \textbf{50} & \\
  \bottomrule
\end{tabular}
\end{center}

% =============================================================================
\section{Materiales y recursos}
% =============================================================================

\begin{itemize}
  \item \textbf{Cuaderno:}
        \texttt{sesiones/sesion19\_proyectos\_multiarchivo.ipynb}
  \item \textbf{Notas alumno:}
        \texttt{notas\_alumno/pdf/sesion19\_proyectos\_multiarchivo.pdf}
  \item \textbf{Material inicial:}
        \texttt{materialInicial/sesion19/} (proyecto mínimo y proyecto1)
\end{itemize}

\end{document}
